\documentclass[11pt]{article}

\usepackage[margin=1in]{geometry}
\usepackage[T1]{fontenc}
\usepackage[utf8]{inputenc}
\usepackage{lmodern}
\usepackage{microtype}
\usepackage{amsmath,amssymb,amsthm}
\usepackage{mathtools}
\usepackage{hyperref}
\usepackage{booktabs}
\usepackage{enumitem}
\usepackage{listings}
\usepackage{xcolor}

\hypersetup{
  colorlinks=true,
  linkcolor=blue!50!black,
  citecolor=blue!50!black,
  urlcolor=blue!50!black
}

\lstdefinelanguage{Lean4}{
  morekeywords={
    import,namespace,end,open,section,noncomputable,abbrev,
    def,theorem,lemma,where,by,fun,let,in,have,show,exact,
    intro,apply,simpa,calc
  },
  sensitive=true,
  morecomment=[l]{--},
  morecomment=[s]{/-}{-/},
  morestring=[b]"
}

\lstset{
  language=Lean4,
  basicstyle=\ttfamily\small,
  keywordstyle=\bfseries\color{blue!50!black},
  commentstyle=\itshape\color{green!40!black},
  stringstyle=\color{red!50!black},
  showstringspaces=false,
  breaklines=true,
  frame=single,
  framerule=0.3pt,
  xleftmargin=0.5em,
  xrightmargin=0.5em,
  keepspaces=true,
  columns=fullflexible,
  literate=
    {¬}{{$\neg$}}1
    {∃}{{$\exists$}}1
    {→}{{$\to$}}1
    {≤}{{$\leq$}}1
}

\newtheorem{theorem}{Theorem}[section]
\newtheorem{lemma}[theorem]{Lemma}
\theoremstyle{remark}
\newtheorem{remark}[theorem]{Remark}

\title{P vs NP Grassroots:\\
Ground-Up Background Theory for the Goertzel Program}
\author{Oruzi / Codex collaboration}
\date{\today}

\begin{document}
\maketitle

\begin{abstract}
This note records the \emph{grassroots} side of the current Lean work around Ben
Goertzel's proposed route to $P \neq NP$.  It is intentionally not the blocker note.
The companion crux-first analysis lives separately in \texttt{pnp\_crux.tex}.

The goal here is narrower and more constructive: build reusable background theory
from the ground up, using mathlib and existing Mettapedia infrastructure wherever
possible, so that if the high-level proof strategy can be repaired, it lands in a
clean machine-checked framework rather than in ad hoc prose.

The current grassroots state now includes explicit encoded hypothesis classes,
finite empirical-risk minimization, finite deceptive-sample bounds, threshold
recovery lemmas, a rational uniform-rate layer, and explicit uniform success
bounds for exact ERM recovery under finite sampling.  It now also includes the
first weighted finite-sampling steps beyond uniform draws: an exact
one-predictor agreement law and a family-level weighted union bound over finite
encoded classes, together with weighted exact-recovery lower bounds for ERM.
This isolates the exact positive shape that a successful switching theorem would
need: not mere locality, but locality together with a concrete compression
theorem and a clean finite learning interface.
\end{abstract}

\tableofcontents

\section{Scope}

The October 2025 Goertzel manuscript
\emph{``P != NP: A Non-Relativizing Proof via Quantale Weakness and Geometric
Complexity''} proposes a route from a structured Unique-SAT ensemble to a
distributional lower bound, and then to a contradiction with the upper bound that
would follow from $P=NP$.

This document does \emph{not} ask where that route fails.  That is the job of the
companion crux note.  Instead it asks a different question:

\begin{quote}
What background theory would we want in Lean even if the final proof still needs
repair?
\end{quote}

That question leads to a program with three design constraints:
\begin{enumerate}[leftmargin=1.5em]
  \item Prefer general infrastructure over one-off rebuttal lemmas.
  \item Reuse mathlib computability and existing Mettapedia semantics rather than
  inventing parallel foundations.
  \item State optimistic bridge theorems in a way that makes missing hypotheses
  explicit, so later fixes can slot into them cleanly.
\end{enumerate}

\section{Existing Lean Base}

Several ingredients already exist in the local formalization and are useful for a
ground-up PNP effort.

\begin{center}
\begin{tabular}{ll}
\toprule
\textbf{Module} & \textbf{Grassroots role} \\
\midrule
\texttt{Mathlib.Computability.Partrec} & partial recursive computation infrastructure \\
\texttt{KolmogorovComplexity/Basic.lean} & universal algorithm and plain complexity \\
\texttt{Hyperseed/Ultrainfinitism.lean} & observer-relative available-region semantics \\
\texttt{Hyperseed/Basic.lean} & closure/cascade interfaces on accessible observations \\
\texttt{PNP/HypothesisClass.lean} & encoded classifier families and code-budget bounds \\
\texttt{PNP/FiniteERM.lean} & empirical-risk minimization kernel for finite encoded families \\
\texttt{PNP/FiniteConsistencyBound.lean} & finite uniform consistency bound for deceptive samples \\
\texttt{PNP/FiniteUniformRecovery.lean} & threshold corollary guaranteeing exact recovery on some sample \\
\texttt{PNP/FiniteUniformRate.lean} & rational uniform-rate bounds for deception and exact recovery \\
\texttt{PNP/FiniteUniformSuccess.lean} & miss-ratio decay and explicit uniform success bounds \\
\texttt{PNP/FiniteIIDAgreement.lean} & exact weighted one-predictor agreement mass under a finite PMF \\
\texttt{PNP/FiniteIIDFamilyBound.lean} & weighted deceptive-mass union bound for finite encoded families \\
\texttt{PNP/FiniteIIDRecovery.lean} & weighted exact ERM recovery lower bounds under finite PMFs \\
\bottomrule
\end{tabular}
\end{center}

The point of the table is simple: the grassroots route is not starting from zero.
We already have a serious computability substrate, an observer-relative semantics
layer, and a new optimistic interface for small hypothesis classes.

\section{Encoded Hypothesis Families}

The key new grassroots abstraction is that a hypothesis class should be represented
by a finite code space together with a decoder.

\begin{quote}
code budget first, realized class second.
\end{quote}

Formally, \texttt{HypothesisClass.lean} introduces an \emph{encoded family}
\[
  \mathrm{decode} : \mathrm{Code} \to (\mathrm{Input} \to \mathrm{Output}),
\]
where \(\mathrm{Code}\) is finite.  The realized hypothesis class is the range of
\(\mathrm{decode}\).

\begin{theorem}[Realized class is bounded by code space]
For any encoded family \(H\),
\[
  |\mathrm{realized}(H)| \le |\mathrm{Code}(H)|.
\]
\end{theorem}

\noindent
This is the theorem \texttt{EncodedFamily.card\_realized\_le} in Lean.

The specialization relevant to the Goertzel program is the bit-coded case.

\begin{theorem}[Bit-budget bound]
If a family of Boolean classifiers is decoded from \(s\) bits, then it realizes at
most \(2^s\) distinct classifiers.
\end{theorem}

\noindent
This is formalized as
\texttt{BitEncodedClassifierFamily.card\_realized\_le}.

This theorem is elementary, but it captures the exact shape of the optimistic
repair that would be needed later.  A switching theorem cannot stop at ``the new
predictor is local.''  It must continue to a statement of the form:

\begin{quote}
the switched predictors lie in a family decoded from \(\mathrm{polylog}(m)\) bits.
\end{quote}

Only then does finite-class learning theory become genuinely available.

\section{Finite-Class ERM Kernel}

The next honest step after a finite encoded family is not yet a PAC theorem.
It is the finite combinatorial kernel of empirical risk minimization.

The new module \texttt{FiniteERM.lean} defines labeled samples as finite lists of
pairs \((x,y)\), empirical error as the number of misclassified examples in the
sample, and exact sample consistency as zero empirical error.

\begin{theorem}[Finite ERM existence]
Let \(H\) be a nonempty encoded family.  For every finite labeled sample \(S\),
there exists a realized predictor \(h^\star \in \mathrm{realized}(H)\) such that
\[
  \mathrm{err}_S(h^\star) \le \mathrm{err}_S(g)
  \qquad
  \text{for all } g \in \mathrm{realized}(H).
\]
\end{theorem}

\noindent
This is formalized as
\texttt{EncodedFamily.exists\_realized\_empiricalRiskMinimizer}.

\begin{theorem}[Zero-error transfer through ERM]
If some code in the family fits the sample exactly, then the chosen ERM
predictor also fits the sample exactly.
\end{theorem}

\noindent
This is formalized as
\texttt{EncodedFamily.empiricalRiskPredictor\_fitsSample\_of\_exists\_code\_fits}.

This is still modest, but it matters.  Combined with the code-budget theorem
from \texttt{HypothesisClass.lean}, it gives the first clean bridge from
``small encoded family'' to an actual learning-theoretic construction.

\section{Finite Uniform Consistency Bound}

The next step now formalized is a fully combinatorial finite-domain bound.  We
do not yet assume an arbitrary distribution or invoke measure theory.  Instead,
we count how many length-\(m\) input samples can make a wrong predictor look
perfectly consistent with the target labels.

\begin{theorem}[Wrong predictors fit only exponentially few samples]
Let \(X\) be a finite input domain.  If \(h \neq f\), then the number of
length-\(m\) samples on which \(h\) agrees with the target \(f\) at every sampled
point is at most
\[
  (|X|-1)^m.
\]
\end{theorem}

\noindent
This is formalized as
\texttt{card\_consistentSamples\_le\_of\_exists\_disagreement} in
\texttt{FiniteConsistencyBound.lean}.

\begin{theorem}[Finite encoded families have few deceptive samples]
Let \(H\) be an encoded family over a finite input domain \(X\).  The number of
length-\(m\) samples on which some wrong code in \(H\) still matches all target
labels is at most
\[
  |\mathrm{Code}(H)| \cdot (|X|-1)^m.
\]
\end{theorem}

\noindent
This is formalized as
\texttt{EncodedFamily.card\_deceptiveSamples\_le}.

\begin{theorem}[ERM recovers the target off the deceptive set]
If the target predictor is realized by the encoded family, then on every
non-deceptive sample the chosen ERM predictor equals the target predictor
exactly.
\end{theorem}

\noindent
This is formalized as
\texttt{EncodedFamily.empiricalRiskPredictor\_eq\_target\_of\_not\_deceptive}.

This is not yet the full finite-class generalization theorem one would state in
terms of \(\log |H|\) and an arbitrary data distribution.  But it is already the
right union-bound skeleton in a completely concrete finite setting.

\section{Finite Uniform Recovery Threshold}

The next corollary is the cleanest possible extraction from the deceptive-sample
bound.  If the deceptive set is strictly smaller than the whole length-\(m\)
sample space, then at least one non-deceptive sample exists.

\begin{theorem}[Threshold for existence of a good sample]
Let \(H\) be an encoded family over a finite input domain \(X\).  If
\[
  |\mathrm{Code}(H)| \cdot (|X|-1)^m < |X|^m,
\]
then there exists a length-\(m\) sample outside the deceptive set.
\end{theorem}

\noindent
This is formalized as
\texttt{EncodedFamily.exists\_nondeceptiveSample\_of\_bound\_lt} in
\texttt{FiniteUniformRecovery.lean}.

\begin{theorem}[Threshold corollary for exact ERM recovery]
If the target predictor is realized by the encoded family and
\[
  |\mathrm{Code}(H)| \cdot (|X|-1)^m < |X|^m,
\]
then there exists a length-\(m\) sample on which the chosen ERM predictor equals
the target predictor exactly.
\end{theorem}

\noindent
This is formalized as
\texttt{EncodedFamily.exists\_sample\_empiricalRiskPredictor\_eq\_target\_of\_bound\_lt}.

This is still finite and non-asymptotic, but it clarifies the next job.  To move
from existence to a real learning theorem, we now need a probability interface
that turns the combinatorial ratio behind the threshold into an explicit success
bound under sampling assumptions.

\section{Finite Uniform Rate Layer}

The next natural refinement is still finite, but no longer purely existential.
Instead of asking only whether some good sample exists, we measure how large the
good and bad regions are inside the finite sample space.

The module \texttt{FiniteUniformRate.lean} defines rational uniform densities on
the point-sample space:
\begin{enumerate}[leftmargin=1.5em]
  \item the deceptive rate,
  \item the non-deceptive rate,
  \item the exact-recovery rate for ERM.
\end{enumerate}

\begin{theorem}[Uniform deceptive-rate bound]
Let \(H\) be an encoded family over a nonempty finite input domain \(X\).  Then
the uniform fraction of length-\(m\) samples that are deceptive is at most
\[
  \frac{|\mathrm{Code}(H)| \cdot (|X|-1)^m}{|X|^m}.
\]
\end{theorem}

\noindent
This is formalized as
\texttt{EncodedFamily.deceptiveRate\_le\_bound} in
\texttt{FiniteUniformRate.lean}.

\begin{theorem}[Exact recovery dominates the non-deceptive region]
If the target predictor is realized by the encoded family, then the uniform exact
recovery rate of the chosen ERM predictor is at least
\[
  1 - \mathrm{rate}_{\mathrm{deceptive}}.
\]
\end{theorem}

\noindent
This is formalized as
\texttt{EncodedFamily.exactRecoveryRate\_ge\_one\_sub\_deceptiveRate}.

\begin{theorem}[Positive exact recovery rate under the threshold]
If the target predictor is realized and
\[
  |\mathrm{Code}(H)| \cdot (|X|-1)^m < |X|^m,
\]
then the exact recovery rate is strictly positive.
\end{theorem}

\noindent
This is formalized as
\texttt{EncodedFamily.exactRecoveryRate\_pos\_of\_bound\_lt}.

This is still not a full PAC-style theorem under arbitrary distributions.  But it
does close the uniform finite-sampling gap between ``some good sample exists'' and
``good samples occupy an explicitly bounded fraction of the sample space.''

\section{Finite Uniform Success Bounds}

The next cleanup step is to express the finite uniform rate bound using the
single-step miss ratio
\[
  \rho_X := \frac{|X|-1}{|X|}.
\]
For any nonempty finite input domain, \(\rho_X < 1\), so the finite deceptive
rate decays exponentially in the sample length \(m\).

The module \texttt{FiniteUniformSuccess.lean} packages the previous rate theorem
into this cleaner interface.

\begin{theorem}[Miss-ratio deceptive bound]
Let \(H\) be an encoded family over a nonempty finite input domain \(X\).  Then
\[
  \mathrm{rate}_{\mathrm{deceptive}}
  \le
  |\mathrm{Code}(H)| \cdot \rho_X^m.
\]
\end{theorem}

\noindent
This is formalized as
\texttt{EncodedFamily.deceptiveRate\_le\_codeMul\_uniformMissRatio\_pow}.

\begin{theorem}[Uniform exact recovery lower bound]
If the target predictor is realized by the encoded family, then
\[
  \mathrm{rate}_{\mathrm{exact\ recovery}}
  \ge
  1 - |\mathrm{Code}(H)| \cdot \rho_X^m.
\]
\end{theorem}

\noindent
This is formalized as
\texttt{EncodedFamily.exactRecoveryRate\_ge\_one\_sub\_codeMul\_uniformMissRatio\_pow}.

\begin{theorem}[\(\varepsilon\)-style corollary]
If
\[
  |\mathrm{Code}(H)| \cdot \rho_X^m \le \varepsilon,
\]
then the exact recovery rate is at least \(1 - \varepsilon\).
\end{theorem}

\noindent
This is formalized as
\texttt{EncodedFamily.exactRecoveryRate\_ge\_one\_sub\_of\_codeMul\_uniformMissRatio\_pow\_le}.

This is still uniform finite-sample learning rather than full distributional
generalization.  But it is the first version of the grassroots story with an
explicit success lower bound of the form ``one minus an exponentially decaying
failure term.''

\section{Finite Weighted Agreement and Family Bounds}

The next honest move is to relax the uniform-sampling assumption while still
staying fully finite.  Instead of sampling uniformly from a finite input domain,
we let one finite probability mass function \(\mu\) weight the inputs.

The module \texttt{FiniteIIDAgreement.lean} handles the one-predictor case.  For
one target \(f\) and one predictor \(h\), it defines:
\begin{enumerate}[leftmargin=1.5em]
  \item the one-step agreement mass of \(h\) with \(f\) under \(\mu\),
  \item the product mass of a length-\(m\) sample under independent draws from \(\mu\),
  \item the total mass of samples on which \(h\) agrees with \(f\) everywhere.
\end{enumerate}

\begin{theorem}[Exact weighted agreement law]
For any finite input type and finite PMF \(\mu\),
\[
  \mathrm{mass}_{\mu}(\text{all-$m$-point samples consistent with } h)
  =
  \bigl(\mathrm{agreementMass}_{\mu}(f,h)\bigr)^m.
\]
\end{theorem}

\noindent
This is formalized as
\texttt{consistentSampleMass\_eq\_agreementMass\_pow}.

\begin{theorem}[Weighted decay corollary]
If the one-step agreement mass is at most \(q\), then the total mass of
length-\(m\) consistent samples is at most \(q^m\).
\end{theorem}

\noindent
This is formalized as
\texttt{consistentSampleMass\_le\_pow\_of\_agreementMass\_le}.

The new module \texttt{FiniteIIDFamilyBound.lean} lifts this to the encoded-family
level.

\begin{theorem}[Weighted deceptive-family union bound]
Let \(H\) be a finite encoded family and \(\mu\) a finite PMF on the input type.
Then the total \(\mu^m\)-mass of deceptive samples is at most
\[
  \sum_{c \in \mathrm{BadCodes}(H,f)}
    \bigl(\mathrm{agreementMass}_{\mu}(f,\mathrm{decode}(c))\bigr)^m.
\]
\end{theorem}

\noindent
This is formalized as
\texttt{EncodedFamily.deceptiveSampleMass\_le\_badCodeAgreementMassSum}.

\begin{theorem}[Uniform weighted family corollary]
If every bad code in the family has one-step agreement mass at most \(q\), then
the total deceptive mass is at most
\[
  |\mathrm{Code}(H)| \cdot q^m.
\]
\end{theorem}

\noindent
This is formalized as
\texttt{EncodedFamily.deceptiveSampleMass\_le\_codeCard\_mul\_pow\_of\_agreementMass\_le}.

This is the first genuinely family-level finite distributional theorem in the
grassroots development.  The next real probabilistic step is no longer a union
bound itself, but a recovery or generalization theorem built on top of this
weighted family bound.

\section{Finite Weighted Recovery Bounds}

The module \texttt{FiniteIIDRecovery.lean} closes the next natural loop.  It
turns the weighted deceptive-mass bound into a weighted success statement for the
chosen ERM predictor.

\begin{theorem}[Weighted exact recovery dominates the nondeceptive mass]
If the target predictor is realized by the encoded family, then the total
\(\mu^m\)-mass of exact ERM recovery is at least
\[
  1 - \mathrm{deceptiveMass}_{\mu}.
\]
\end{theorem}

\noindent
This is formalized as
\texttt{EncodedFamily.exactRecoverySampleMass\_ge\_one\_sub\_deceptiveSampleMass}.

\begin{theorem}[Weighted exact recovery from bad-code agreement masses]
If the target predictor is realized by the encoded family, then the exact
recovery mass is at least
\[
  1 - \sum_{c \in \mathrm{BadCodes}(H,f)}
    \bigl(\mathrm{agreementMass}_{\mu}(f,\mathrm{decode}(c))\bigr)^m.
\]
\end{theorem}

\noindent
This is formalized as
\texttt{EncodedFamily.exactRecoverySampleMass\_ge\_one\_sub\_badCodeAgreementMassSum}.

\begin{theorem}[Uniform weighted exact recovery corollary]
If every bad code has one-step agreement mass at most \(q\), then the exact
recovery mass is at least
\[
  1 - |\mathrm{Code}(H)| \cdot q^m.
\]
\end{theorem}

\noindent
This is formalized as
\texttt{EncodedFamily.exactRecoverySampleMass\_ge\_one\_sub\_codeCard\_mul\_pow\_of\_agreementMass\_le}.

At this point the grassroots development has reached a fully finite,
distribution-weighted recovery theorem for ERM over encoded families.  The next
step is to lift this from exact recovery on a finite sample space to a genuine
finite-class generalization statement.

\section{Hyperseed as a Local-View Surface}

The existing Hyperseed formalization already provides a clean observer-relative
semantics.

Given a perspective \(P\), a budget \(B\), and a guard region \(G\), it defines
an \emph{available region}
\[
  \mathrm{availableRegion}(P,B,G)
  =
  \mathrm{observableUniverse}(P)
  \cap
  \mathrm{nearEurycosm}(P,B)
  \cap
  G.
\]

For the grassroots PNP program, this matters because it gives a principled way to
talk about what a decoder can even access.  Instead of treating locality as an
informal phrase, we can work with a typed input surface:
\[
  \mathrm{AvailableInput}(P,B,G)
  := \{x : \mathrm{World} \mid x \in \mathrm{availableRegion}(P,B,G)\}.
\]

The new Lean module proves the corresponding encoded-family bound on such an
available input surface:

\begin{theorem}[Available-region code-budget bound]
An \(s\)-bit classifier family over one Hyperseed available region still realizes
at most \(2^s\) classifiers.
\end{theorem}

\noindent
This is formalized as
\texttt{card\_realized\_availableFamily\_le\_of\_bitBudget}.

The theorem is intentionally modest.  It does not claim that the Goertzel local
view \emph{is} a Hyperseed perspective.  It says that Hyperseed already gives us a
good general language for bounded observational access, and that this language can
interoperate with the encoded-hypothesis-class infrastructure.

\section{Kolmogorov Foundations Already Present}

The local repository also already contains a substantial Kolmogorov-complexity seed
under \texttt{Mettapedia/Computability/KolmogorovComplexity/}.

In particular, \texttt{Basic.lean} defines:
\begin{enumerate}[leftmargin=1.5em]
  \item binary strings as \texttt{List Bool},
  \item partial algorithms,
  \item plain Kolmogorov complexity \(C[U](x)\),
  \item a universal algorithm built from mathlib's partial-recursive code,
  \item the usual invariance theorem for universal algorithms.
\end{enumerate}

This is not yet the time-bounded conditional complexity \(K_{\mathrm{poly}}\)
that the Goertzel program uses.  But it matters for the grassroots route because it
means the eventual \(K_{\mathrm{poly}}\) development should be an extension of an
existing computability layer rather than a fresh bespoke definition dropped into a
late-stage contradiction proof.

In practical terms, the next computability milestone is:

\begin{quote}
define conditional and polynomial-time-bounded complexity on top of the existing
universal-machine layer, instead of axiomatizing them only at the paper level.
\end{quote}

\section{What the Grassroots Program Needs Next}

The current background theory suggests a fairly crisp optimistic roadmap.

\subsection{Finite-Class Learning Layer}

The ERM existence layer, finite uniform consistency bound, threshold recovery
corollary, rational uniform-rate layer, miss-ratio success bound, and weighted
one-predictor and finite-family weighted bounds are now in place.  The next
finite-class learning
steps are:
\begin{enumerate}[leftmargin=1.5em]
  \item a generalization bound stated in terms of \(\log |H|\),
  \item a probability interface for i.i.d. or distributional assumptions at the family level,
  \item concentration-style control beyond the exact finite-space formulas,
  \item a clean interface separating statistical assumptions from coding facts.
\end{enumerate}

In other words, the current Lean state already handles:
\begin{enumerate}[leftmargin=1.5em]
  \item empirical error on finite samples,
  \item existence of ERM predictors in nonempty finite encoded families,
  \item exact-fit transfer when some code in the family is sample-consistent,
  \item a finite uniform bound on the deceptive sample set,
  \item exact ERM recovery of the target off that deceptive set,
  \item a strict-cardinality threshold guaranteeing the existence of a good sample,
  \item a rational upper bound on the deceptive sample rate,
  \item a rational lower bound on the exact recovery rate,
  \item strict positivity of that exact recovery rate under the threshold,
  \item an explicit miss-ratio success bound of the form \(1 - |H|\rho_X^m\),
  \item an exact weighted one-predictor law of the form \(\mathrm{mass} = \alpha^m\),
  \item a weighted family union bound of the form
  \(\mathrm{deceptiveMass} \le \sum_{h \in H_{\mathrm{bad}}}\alpha_h^m\),
  \item a weighted exact recovery bound of the form
  \(\mathrm{recoveryMass} \ge 1 - \sum_{h \in H_{\mathrm{bad}}}\alpha_h^m\).
\end{enumerate}

\noindent
What remains is to lift that finite uniform rate layer into a real finite-class
generalization theorem.

\subsection{Concrete Decoder Subclasses}

The encoded-family theorem is only an interface.  It still needs a nontrivial
source of small classes.  The likely optimistic candidates are concrete local
decoders on tree-like neighborhoods:
\begin{enumerate}[leftmargin=1.5em]
  \item belief-propagation style update rules,
  \item GF(2)-linear constraint summaries from the VV layer,
  \item bounded-depth compositions of those primitives.
\end{enumerate}

If those can be encoded uniformly in \(\mathrm{polylog}(m)\) bits, then the
switching theorem has a plausible landing zone.

\subsection{Bridging to the Full PNP Route}

Only after those layers are built does it make sense to reconnect to the full
proof program:
\begin{enumerate}[leftmargin=1.5em]
  \item define the distributional ensemble cleanly,
  \item define local views as typed accessible surfaces,
  \item prove the switched family lies in a small encoded subclass,
  \item connect success bounds to a properly formalized \(K_{\mathrm{poly}}\),
  \item finally compare with the upper bound under \(P=NP\).
\end{enumerate}

\section{Current Status}

At present, the grassroots side has several real machine-checked gains and one
major pre-existing base layer:

\begin{center}
\begin{tabular}{lll}
\toprule
\textbf{Layer} & \textbf{Status} & \textbf{Artifact} \\
\midrule
encoded hypothesis families & proved & \texttt{HypothesisClass.lean} \\
finite-class ERM kernel & proved & \texttt{FiniteERM.lean} \\
finite uniform consistency bound & proved & \texttt{FiniteConsistencyBound.lean} \\
finite uniform recovery threshold & proved & \texttt{FiniteUniformRecovery.lean} \\
finite uniform rate layer & proved & \texttt{FiniteUniformRate.lean} \\
finite uniform success bounds & proved & \texttt{FiniteUniformSuccess.lean} \\
finite weighted one-predictor agreement law & proved & \texttt{FiniteIIDAgreement.lean} \\
finite weighted family union bound & proved & \texttt{FiniteIIDFamilyBound.lean} \\
finite weighted recovery bounds & proved & \texttt{FiniteIIDRecovery.lean} \\
Hyperseed available-region reuse & proved & \texttt{HypothesisClass.lean} + \texttt{Hyperseed} \\
plain Kolmogorov infrastructure & already present & \texttt{KolmogorovComplexity/Basic.lean} \\
finite-class generalization bounds & not yet formalized & next target \\
\(K_{\mathrm{poly}}\) layer & not yet formalized & later target \\
\bottomrule
\end{tabular}
\end{center}

This is not the final proof.  It is the right sort of groundwork.

\section{Conclusion}

The grassroots program should stay separate from the crux program.

The crux note asks where the present paper overreaches.  This note asks what
general machinery we would still want even if Ben patches the overreach later.
On that constructive side, the current Lean result is that small encoded
hypothesis families, finite empirical-risk minimizers, finite deceptive-sample
bounds, finite uniform recovery thresholds, rational exact-recovery rates,
explicit miss-ratio success bounds, weighted one-predictor agreement laws,
weighted family deceptive-mass bounds, weighted exact recovery bounds, and
observer-relative input surfaces can already be handled cleanly inside one
coherent Mettapedia/mathlib formal background layer.

That is enough to sharpen the next optimistic task.  We should not try to prove
\(P \neq NP\) in one jump.  We should first prove that the intended switched
predictors live in a genuinely small encoded class, and only then lift the rest
of the argument on top of that foundation.

\end{document}
